\chapter{Objetivos do Curso}

O curso de \ppccurso destina-se a preparar engenheiros(as) para atuar de forma integrada nas áreas de hardware e software, bem como em suas interações e aplicações sistêmicas. O profissional formado pela FEELT/UFU será capaz de pensar de forma holística, agir com base em conhecimentos técnico-científicos sólidos e manter postura crítica, ética e inovadora. Espera-se que demonstre iniciativa, competência social, responsabilidade profissional e compromisso com o desenvolvimento sustentável e tecnológico da sociedade.

\section{Objetivo Geral}

Formar profissionais generalistas na área da \ppccurso, com domínio técnico-científico, capacidade de absorver e desenvolver novas tecnologias, e aptidão para atuar de forma crítica e criativa na identificação e resolução de problemas. O(a) egresso(a) deverá considerar os aspectos políticos, econômicos, sociais, ambientais e culturais da prática profissional, mantendo uma postura ética e humanística em consonância com as transformações tecnológicas globais e as demandas da sociedade brasileira.

\section{Objetivos Específicos}

Considerando as Diretrizes Curriculares Nacionais~\cite{MEC:CNE:CES:5:2016} e os princípios formativos contemporâneos do \textit{Computing Curricula 2020}~\cite{ACM:CC2020}, o curso de \ppccurso busca:

\begin{itemize}
    \item Estimular a postura investigativa, a mentalidade inovadora e a produção de conhecimento articulada ao domínio de fundamentos tecnológicos;
    \item Promover a integração curricular entre as formações geral, profissional e específica, assegurando progressão adequada e interdisciplinaridade;
    \item Capacitar os(as) futuros(as) engenheiros(as) da computação a incorporar os valores e objetivos do desenvolvimento sustentável em sua prática profissional, utilizando a engenharia como ferramenta para responder aos desafios globais, conforme preconizado pelo EESD;
    \item Garantir atualização permanente do currículo, flexibilizando os conteúdos mais sujeitos à obsolescência por meio de disciplinas optativas e módulos eletivos de aprofundamento;
    \item Favorecer a realização de atividades interdisciplinares integradas, visando o diálogo entre diferentes áreas de conhecimento no mesmo período letivo;
    \item Integrar teoria e prática, conectando conteúdos curriculares às realidades da atuação profissional, por meio de projetos, oficinas, estudos de caso e experiências aplicadas;
    \item Reconhecer e incorporar ao percurso formativo as competências desenvolvidas em contextos não formais de aprendizagem (iniciação científica, monitorias, estágios, extensão, vivência profissional);
    \item Acompanhar as transformações do mercado de trabalho e da sociedade, fomentando respostas éticas, criativas e socialmente responsáveis às demandas consolidadas e emergentes;
    \item Assegurar a formação de profissionais capazes de:
    \begin{itemize}
        \item Compreender e atuar sobre questões sociais, éticas, legais, políticas e humanísticas que permeiam a prática profissional;
        \item Antecipar e atender de forma estratégica as necessidades da sociedade por meio da aplicação crítica da computação e suas tecnologias;
        \item Contribuir com visão criativa e transformadora para a evolução da engenharia e do conhecimento técnico-científico;
        \item Empreender, liderar e cooperar em equipes multidisciplinares, atuando em cenários locais e globais com senso de responsabilidade;
        \item Utilizar racional e transdisciplinarmente os recursos disponíveis, com foco na sustentabilidade e na inovação.
    \end{itemize}
\end{itemize}